% ID: SEL-D0052-OB-E02
% ESTADO: validada. S053.
\textbf{a)} Las dos ramas polinómicas cubren $\mathbb R$. En la unión, $\lim_{x\to1^-}f(x)=0=f(1)=\lim_{x\to1^+}f(x)$, de modo que $f$ es continua en todo su dominio.

\textbf{b)} En $x<1$, $f'(x)=3x^2\ge0$ y la cúbica es estrictamente creciente, incluso al pasar por $0$. Para $x>1$, $f'(x)=4-2x$, positiva hasta $2$ y negativa después. En $1$ las pendientes laterales $3$ y $2$ difieren, pero el crecimiento continúa. Hay un único máximo relativo y absoluto en $x=2$, con $f(2)=\boxed{1}$. No hay mínimo, pues la función decrece sin cota inferior cuando $x\to+\infty$.

\textbf{c)} A la izquierda de $1$, $f''(x)=6x$. A la derecha, $f''(x)=-2$. Por tanto, la gráfica es cóncava hacia abajo en $(-\infty,0)$ y $(1,\infty)$ y hacia arriba en $(0,1)$. En $(0,-1)$ hay una inflexión con tangente horizontal. En $(1,0)$ también cambia la curvatura, pero la unión es angulosa y carece de tangente única. Si se exige tangente para llamar inflexión a un punto, sólo se incluye $(0,-1)$.
